\section{Evaluation}
\label{sec:mllm_eval}
Evaluation is an essential part of developing MLLMs since it provides feedback for model optimization and helps to compare the performance of different models. Compared with evaluation methods of traditional multimodal models, the evaluation of MLLMs exhibits several new traits: (1) Since MLLMs are generally versatile, it is important to evaluate MLLMs comprehensively. (2) MLLMs exhibit many emergent capabilities that require special attention (\eg OCR-free math reasoning) and thus require new evaluation schemes.
The evaluation of MLLMs can be broadly categorized into two types according to the question genres, including closed-set and open-set. 

\subsection{Closed-set}
Closed-set questions refer to a type of question where the possible answer options are predefined and limited to a finite set. 
The evaluation is usually performed on task-specific datasets. 
In this case, the responses can be naturally judged by benchmark metrics~\cite{instructblip, llava, x-llm, multiinstruct, llama-adapter, llama-adapter-v2, chatbridge, lavin}. 
For example, InstructBLIP~\cite{instructblip} reports the accuracy on ScienceQA~\cite{scienceqa}, as well as the CIDEr score~\cite{vedantam2015cider} on NoCaps~\cite{agrawal2019nocaps} and Flickr30K~\cite{young2014image}. 
The evaluation settings are typically zero-shot~\cite{instructblip, multiinstruct, chatbridge, m3it} or finetuning~\cite{instructblip, llava, x-llm, llama-adapter, llama-adapter-v2, lavin, llava-med, m3it}. 
The first setting often selects a wide range of datasets covering different general tasks and splits them into held-in and held-out datasets. 
After tuning on the former, zero-shot performance is evaluated on the latter with unseen datasets or even unseen tasks. 
In contrast, the second setting is often observed in the evaluation of domain-specific tasks. 
For example, LLaVA~\cite{llava} and LLaMA-Adapter~\cite{llama-adapter} report finetuned performance on ScienceQA~\cite{scienceqa}. 
LLaVA-Med~\cite{llava-med} reports results on biomedical VQA~\cite{he2020pathvqa, lau2018dataset, liu2021slake}.

The above evaluation methods are usually limited to a small range of selected tasks or datasets, lacking a comprehensive quantitative comparison. 
To this end, some efforts have endeavored to develop new benchmarks specially designed for MLLMs~\cite{mme,mmbench,mm-vet,seed-bench,mathvista,mmmu,hallusionbench}. 
For example, Fu \etal~\cite{mme} construct a comprehensive evaluation benchmark MME that includes a total of 14 perception and cognition tasks.
All instruction-answer pairs in MME are manually designed to avoid data leakage.
MMBench~\cite{mmbench} is a benchmark specifically designed for evaluating multiple dimensions of model capabilities, using ChatGPT to match open responses with pre-defined choices.
Video-ChatGPT~\cite{video-chatgpt} and Video-Bench~\cite{ning2023video} focus on video domains and propose specialized benchmarks as well as evaluation tools for assessment.
There are also evaluation strategies designed to evaluate a specific aspect of the model~\cite{multiinstruct}, as exemplified by POPE~\cite{li2023evaluating} for assessment of hallucination degree.


\subsection{Open-set}
In contrast to the closed-set questions, the responses to open-set questions can be more flexible, where MLLMs usually play a chatbot role.
Because the content of the chat can be arbitrary, it would be trickier to judge than the closed-ended output.
The criterion can be classified into manual scoring, GPT scoring, and case study. 
Manual scoring requires humans to assess the generated responses. 
This kind of approach often involves hand-crafted questions that are designed to assess specific dimensions. 
For example, mPLUG-Owl~\cite{mplug-owl} collects a visually related evaluation set to judge capabilities like natural image understanding, diagram, and flowchart understanding.  
Similarly, GPT4Tools~\cite{gpt4tools} builds two sets for the finetuning and zero-shot performance, respectively, and evaluates the responses in terms of thought, action, arguments, and the whole.

Since manual assessment is labor intensive, some researchers have explored rating with GPT, namely GPT scoring. 
This approach is often used to evaluate performance on multimodal dialogue. 
LLaVA~\cite{llava} proposes to score the responses via text-only GPT-4 in terms of different aspects, such as helpfulness and accuracy. 
Specifically, 30 images are sampled from the COCO~\cite{COCO} validation set, each associated with a short question, a detailed question, and a complex reasoning question via self-instruction on GPT-4. 
The answers generated by both the model and GPT-4 are sent to GPT-4 for comparison. 
Subsequent works follow this idea and prompt ChatGPT~\cite{mplug-owl} or GPT-4~\cite{x-llm, chatbridge, lavin, llava-med, m3it} to rate results~\cite{mplug-owl, x-llm, chatbridge, lavin, llava-med} or judge which one is better~\cite{llama-adapter-v2}.


A main issue of applying text-only GPT-4 as an evaluator is that the judge is only based on image-related text content, such as captions or bounding box coordinates, without accessing the image~\cite{llava-med}. 
Thus, it may be questionable to set GPT-4 as the performance upper bound in this case. With the release of the vision interface of GPT, some works~\cite{yin2023woodpecker, li2024red} exploit a more advanced GPT-4V model to assess the performance of MLLMs. For example, Woodpecker~\cite{yin2023woodpecker} adopts GPT-4V to judge the response quality of model answers based on the image. The evaluation is expected to be more accurate than using text-only GPT-4 since GPT-4V has direct access to the image.

A supplementary approach is to compare the different capabilities of MLLMs through case studies.
For instance, some studies evaluate two typical advanced commercial-use models, GPT-4V and Gemini. 
Yang~\etal~\cite{yang2023dawn} perform in-depth qualitative analysis on GPT-4V by crafting a series of samples across various domains and tasks, spanning from preliminary skills, such as caption and object counting, to complex tasks that require world knowledge and reasoning, such as joke understanding and indoor navigation as an embodied agent. 
Wen~\etal~\cite{wen2023road} make a more focused evaluation of GPT-4V by designing samples targeting automatic driving scenarios.
Fu~\etal~\cite{fu2023challenger} carry out a comprehensive evaluation on Gemini-Pro by comparing the model against GPT-4V. The results suggest that GPT-4V and Gemini exhibit comparable visual reasoning abilities in spite of different response styles.
